\documentclass{jsarticle}
\usepackage{amsmath}
\usepackage{amssymb}
\begin{document}
\title{$ AdS_5$多様体と別次元の$AdS_5$\\
素数と素粒子方程式}
\author{Masaaki Yamaguchi}
\date{}
\maketitle

   $$ e^{-2\pi ||psi|}[\eta_{\mu\nu} + \hbar{x}]dx^{\mu}dx^{\nu} 
  +  T^2 d^2 \psi $$
  $$ = e^{-f} + e^f = (u + d) + c, (w + s) + b = -e^{-f} + e^f $$
  この式は、時空にどれだけの原子が凝縮しているかを表していて、
  その上に、この逆数は、密度エネルギーをも示している。
  空間の体積にに原子を商代数にすると、濃度にもなる。
   逆数は、$\text{a=原子のエネルギー}$が全空間を1とすると
  $\text{$1 \over {y \over 
  x}$}$すると、$x \over y $は密度になる。　
  $\text{{x=原子},y=全空間、個数は}$$\text{$y \over x$}$ 
  $$ \square = \text{原子} \to {{[\eta_{\mu\nu} + \hbar(x)]dx^{\mu}dx^{\nu}
  } \over {e^{2 \pi T ||\psi||} }}$$
  $$ ||ds^2|| = e^{2 \pi T||\psi||}\left([\eta_{\mu\nu} + \hbar(x)]d^{\mu}dx^
  {\nu}\right)^{-1} + \left(T^2 d^2 \psi \right)^{-1} $$
  $$||ds^2|| = e^{-2 \pi T||\psi||}[\eta_{\mu\nu} + \hbar(x)]d^{\mu}dx^
  {\nu}] + T^2 d^2 \psi  $$
  $$ (u + d) + c = e^{-f} + e^f, (w + s) + b = e^{-f} + e^{f} $$
  $$ R_{ij}^{'} = - R_{ij} $$

  この$ AdS_5 $多様体の $e^{-2 \pi T ||\psi||} $は原子間距離を表していて、
  $ [\eta + \bar{h}(x)]d^{\mu}dx^{\nu} $で宇宙のD-braneの構造を表している。
  この数式が$T^2 d^2 \psi $とアーベル多様体が包み込んでいる。
  原子が回転するのと、ニュートンリングが回転する宇宙は同じ速さで回転する。　

  $$||ds^2|| = \square + \rho , ||ds^2|| = \nabla \Psi^2 + \psi $$
  $$ \eta_{\mu\nu} + \bar{h}(x) $$
  $$ \text{ploximetry} = \Delta x \Delta p - \Delta p \Delta x + 
  \delta(p,x)  $$
  $$ ||ds^2|| = e^{-2 \pi T ||\psi||} + [\eta_{\mu\nu}  + \bar{h}(x)]
  dx^{\mu\nu} + T^2 d^2 \psi $$
  $$ {d \over df}F = e^f + e^{-f}, {d \over df}\int C dx_m = e^f - e^{-f} $$
  $$ R_{ij}^{'} = - R_{ij} $$
  宇宙と異次元では０だが宇宙だけだと誤差が生じる。これが不確定性原理であり、
  異次元が誤差になり、宇宙と加群すると０になる。
  　$ \delta(p,x) $が隠れた変数となり、異次元で誤差をこの不確定性原理
  と表している。
  全ては素粒子方程式から生成される式達である。
\end{document}
